\documentclass{article}

\usepackage{amsmath, amsthm, amssymb, amsfonts}
\usepackage{thmtools}
\usepackage{graphicx}
\usepackage{setspace}
\usepackage[margin=1.5in]{geometry}
\usepackage{float}
\usepackage{hyperref}
\usepackage[utf8]{inputenc}
\usepackage[english]{babel}
\usepackage{framed}
\usepackage[dvipsnames]{xcolor}
\usepackage{tcolorbox}
\usepackage{verbatim}
\usepackage{txfonts}
\usepackage{qtree}

\newcommand{\R}{\mathbb{R}}
\newcommand{\N}{\mathbb{N}}
\newcommand{\Z}{\mathbb{Z}}
\newcommand{\Q}{\mathbb{Q}}
\newcommand{\PP}{\mathbb{P}}
\newcommand{\E}{\mathbb{E}}
\newcommand{\C}{\mathbb{C}}
\newcommand{\T}{\mathbb{T}}

\setlength{\parindent}{0pt}

\colorlet{LightGray}{White!90!Periwinkle}
\colorlet{LightOrange}{Orange!15}
\colorlet{LightGreen}{Green!15}

\newcommand{\HRule}[1]{\rule{\linewidth}{#1}}

\declaretheoremstyle[name=Theorem,]{thmsty}
\declaretheorem[style=thmsty,numberwithin=section]{theorem}
\tcolorboxenvironment{theorem}{colback=LightGray}

\declaretheoremstyle[name=Example,]{prosty}
\declaretheorem[style=prosty,numberlike=theorem]{example}
\tcolorboxenvironment{example}{colback=LightOrange}

\declaretheoremstyle[name=Definition,]{prcpsty}
\declaretheorem[style=prcpsty,numberlike=theorem]{definition}
\tcolorboxenvironment{definition}{colback=LightGreen}

\setstretch{1.2}
\geometry{
  textheight=9in,
  textwidth=5.5in,
  top=1in,
  headheight=12pt,
  headsep=25pt,
  footskip=30pt
}

% ------------------------------------------------------------------------------

\begin{document}

% ------------------------------------------------------------------------------
% Cover Page and ToC
\title{ \normalsize \textsc{}
  \\ [2.0cm]
  \HRule{1.5pt} \\
  \LARGE \textbf{\uppercase{CAS MA394 Applied Abstract Algebra}
  \HRule{2.0pt} \\ [0.6cm] \LARGE{Spring 2026} \vspace*{10\baselineskip}}
}
\date{}
\author{\textbf{Frank Yang} \\
  Professor Tom Enkosky \\
TTh 2:00 -- 3:15}

\maketitle
\newpage

\tableofcontents
\newpage

\section{Lecture 1 -- 1/20}
\subsection{1.2 Sets and Equivalence Relations}
\begin{definition}
  A \textbf{set} is a well-defined collection of objects; that is, it is defined in such a manner that we can determine for any given object whether or not it belonsgs to the set. The objects that belond to a set are called its \textbf{elements} or \textbf{members}. We denote them with capital letters such as $A$ or $X$.

  \begin{itemize}
    \item If $a$ is an element of a set $A$, we write $a\in A$.
    \item The set without any elements its the \textbf{empty set}, denoted $\emptyset$ or $\{\}$.
    \item If $A$ and $B$ are sets and every element of $A$ is in $B$, then $A$ is a \textbf{subset} of $B$, denoted $A\subseteq B$.
    \item If $A\subset B$ and $B\subset A$, then $A=B$.
    \item If $A\subset B$ and $A\neq B$, then $A$ is a \textbf{proper subset} of $B$, denoted $A\subset B$.
  \end{itemize}
  Some standard sets include:
  \begin{itemize}
    \item $\N = \{1, 2, 3, \ldots\}$, the set of natural numbers (positive integers).
    \item $\Z = \{\ldots, -2, -1, 0, 1, 2, \ldots\}$, the set of integers.
    \item $\Q = \left\{\frac{a}{b} \mid a, b \in \Z, b \neq 0\right\}$, the set of rational numbers.
    \item $\R=(-\infty, \infty)$, the set of real numbers.
    \item $\C=\{a+bi \mid a,b\in \R, i^2=-1\}$, the set of complex numbers.
  \end{itemize}
\end{definition}
We can also build new sets from old sets.
\begin{itemize}
  \item The \textbf{union} of two sets $A$ and $B$ is the set of elements that are in $A$ or $B$ (or both), denoted $A\cup B$.
    \[A\cup B=\{x:x\in A\text{ or } x\in B\}.\]
  \item The \textbf{intersection} of two sets $A$ and $B$ is the set of elements that are in both $A$ and $B$, denoted $A\cap B$.
    \[A\cap B=\{x:x\in A\text{ and } x\in B\}.\]
  \item The \textbf{complement} of a set $A$ (with respect to a universal set $U$) is the set of elements in $U$ that are not in $A$, denoted $A',A^c,\bar{A}$.
    \[A'=\{x:x\in U\text{ and } x\notin A\}.\]
  \item The \textbf{difference} of two sets $A$ and $B$ is the set of elements that are in $A$ but not in $B$, denoted $A\setminus B$ or $A-B$.
    \[A\setminus B=A\cap B'=\{x:x\in A\text{ and } x\notin B\}.\]
\end{itemize}
\begin{definition}
  \textbf{De Morgan's Laws} state that for any two sets $A$ and $B$,
  \begin{align*}
    (A\cup B)' &= A' \cap B' \\
    (A\cap B)' &= A' \cup B'.
  \end{align*}
\end{definition}
\begin{definition}
  The \textbf{Cartesian product} of two sets $A$ and $B$ is the set of ordered pairs $(a,b)$ where $a\in A$ and $b\in B$, denoted $A\times B$.
  \[A\times B=\{(a,b):a\in A, b\in B\}.\]
\end{definition}
An example of a Cartesian product is $\R^2=\R\times \R=\{(x,y):x,y\in \R\}$, the set of all points in the plane.
\begin{definition}
  A \textbf{relation} $R$ from a set $A$ to a set $B$ is a subset of the Cartesian product $A\times B$.
\end{definition}
\begin{example}
  Let $A=\{0, 1, 2, 3\}, B=\{0, 1, 2, 3\}$. Then $A\times B=\{(a,b):a,b\in A\}$. If we say that $R$ is the set $\{(0, 1), (0, 2), (0, 3), (1, 2), (1, 3), (2, 3)\}\subset A\times B$, then $R$ is a relation from $A$ to $B$.The relation $R$ can be interpreted as "is less than", since for each $(a,b)\in R$, we have $a<b$.
\end{example}
\begin{definition}
  We define a mapping or function $f\subset A\times B$ from a set $A$ to a set $B$ as a special type of relation where each $a\in A$ has a unique $b\in B$ such that $(a,b)\in f$. We denote this by $f:A\to B$ or $A\xrightarrow{f}B$.
  The set $A$ is called the \textbf{domain} of $f$ and $f(A)=\{f(a):a\in A\}\subset B$ is called the range or \textbf{image} of $f$.
\end{definition}
\section{Lecture 2 -- 1/21}
\begin{definition}
  If $f:A\to B$ is a map and the image of $f$ is $B$, i.e., $f(A)=B$, then $f$ is called \textbf{onto} or \textbf{surjective}. In other words, if there exists an $a\in A$ for every $b\in B$ such that $f(a)=b$, then $f$ is onto.
\end{definition}
\begin{definition}
  A map is $\textbf{one-to-one}$ or \textbf{injective} if $a_1\neq a_2$ implies $f(a_1)\neq f(a_2).$ Equivalently, a function is one-to-one if $f(a_1)=f(a_2)$ implies $a_1=a_2$.
\end{definition}
\begin{definition}
  A map $f:A\to B$ is a \textbf{bijection} if it is both one-to-one and onto.
\end{definition}
\begin{definition}
  Given two functions, we can construct a new function by using the range of the first function as the domain of the second function. Let $f:A\to B$ and $g:B\to C$ be mappings. Define a new map $g\circ f:A\to C$ called the \textbf{composition} of $f$ and $g$ from $A$ to $C$, by $g\circ f(x)=g(f(x))$.
\end{definition}
\begin{itemize}
  \item If $S$ is any set, we use $id_S$ or $id$ to denote the \textbf{identity mapping} from $S$ to itself, or $id(s)=s$ for all $s\in S$.
  \item A map $g:B\to A$ is an \textbf{inverse mapping} of$f:A\to B$ if $g\circ f=id_A$ and $f\circ g=id_B$; in other words, the inverse function of a function simply "undoes" the function.
  \item A map is said to be \textbf{invertible} if it has an inverse. We write $f^{-1}$ to denote the inverse of $f$.
\end{itemize}

\section{Lecture 3 -- 1/22}
\begin{definition}
  An \textbf{equivalence relation} on a set $X$ is a relation $R\subset X \times X$ such that
  \begin{itemize}
    \item $(x,x)\in R$ for all $x\in X$ (reflexive property),
    \item $(x,y)\in R$ implies $(y,x)\in R$ (symmetric property),
    \item $(x,y)\in R$ and $(y,z)\in R$ imply $(x,z)\in R$ (transitive property).
  \end{itemize}
  Given an equivalence relation $R$ on a set $X$, we usually write $x\sim y$ instead of $(x,y)\in\R$. If the equivalence relation already has an associated notation such as $=, \cong, \equiv$, we use that notation instead.
\end{definition}
\begin{example}
  Let $X=\{1, 2, 3, 4\}$. A set $R$ that is reflexive, symmetric, and transitive (and thus an equivalence relation) is $R=\{(1,1), (2,2), (3,3), (4,4), (1,2), (2,1), (1, 3), (3, 1), (2, 3), (3, 2)\}$.
\end{example}
\begin{definition}
  A \textbf{partition} $P$ of a set $X$ is a collection of nonempty sets $X_1, X_2, \ldots, X_n$ such that $X_i\cap X_j=\emptyset$ and $\bigcup_{i} X_i = X$. Let $\sim$ be an equivalence relation on a set $X$ and let $x\in X$. Then $[x]=\{y\in X:x\sim y\}$ is called the \textbf{equivalence class} of $x$.
\end{definition}
\begin{example}
  Let $\Z$ be partitioned into $A=\{\text{even integers}\}$ and $B=\{\text{odd integers}\}$.

  $A\cap B=\emptyset$ and $A\cup B-\Z\implies A,B\text{ partition the integers}$. Define a relation $\sim$ on $\Z$ where $x\sim y\iff x,y$ both even or both odd.

  Every $a\in A$ is equivalent to every other $a'\in A$ and every $b\in B$ is equivalent to every other $b'\in B$.

\end{example}

\section{Lecture 4 -- 1/27}
\subsection{2.1 Mathematical Induction}
\begin{definition}
  \textbf{First Principle of Mathematical Induction}: Let $S(n)$ be a statement about integers for $n\in \N$ and suppose $S(n_0)$ is true for some integer $n_0\in\N$. If for all integers $k$ with $k\geq n_0$, $S(k)$ implies that $S(k+1)$ is true, then $S(n)$ is true for all integers $n\geq n_0$.
\end{definition}
\begin{example}
  For all $n\in\N,1+2+3+\cdots+n=\frac{n(n+1)}{2}$.

  The base case is $n=1$:, then $1=\frac{1(1+1)}{2}=1$, so the base case holds.\\
  Inductive hypothesis: Assume $1+2+3+\cdots+n=\frac{n(n+1)}{2}$ for some $n$.\\
  Inductive step: We want to show that the statement is true for $n+1$. Consider $1+2+3+\cdots+n+(n+1)=\frac{n(n+1)}{2}+(n+1)=(n+1)(\frac{n}{2}+1)=(n+1)(\frac{n+2}{2})=\frac{(n+1)((n+1)+1)}{2}$.
\end{example}
\begin{example}
  Prove that
  \[\frac{1}{2}_\frac{1}{6}+\cdots+\frac{1}{n(n+1)}=\frac{n}{n+1}\]
  for $n\in\N$.

  Base: $n=1$, then $\frac{1}{1\cdot 2}=\frac{1}{2}$, so the base case holds.\\
  Inductive hypothesis: Assume $\frac{1}{2}+\frac{1}{6}+\cdots+\frac{1}{n(n+1)}=\frac{n}{n+1}$ for some $n$.\\
  Inductive step: We want to show that the statement is true for $n+1$. Consider
  \begin{align*}
    \frac{1}{2}+\frac{1}{6}+\cdots+\frac{1}{n(n+1)}+\frac{1}{(n+1)(n+2)}&=\frac{n}{n+1}+\frac{1}{(n+1)(n+2)}\\&=\frac{n(n+2)+1}{(n+1)(n+2)}\\&=\frac{(n+1)^2}{(n+1)(n+2)}\\&=\frac{n+1}{n+2}.
  \end{align*}
  Therefore the statement is true for all $n$.
\end{example}

\subsection{2.2 The Division Algorithm}
\begin{definition}
  Let $a$ and $b$ be integers with $b>0$. Then there exist unique integers $q$ and $r$ such that $a=bq+r$ where $0\leq r<b$.

  \begin{itemize}
    \item Let $a$ and $b$ be integers. If $b=ak$ for some integer $k$, we write $a|b$.
    \item An integer $d$ is called a \textbf{common divisor} of $a$ and $b$ if $d|a$ and $d|b$.
    \item The \textbf{greatest common divisor} of $a$ and $b$ is a positive integer $d$ such that $d$ is a common divisor of $a$ and $b$, and if $d'$ is any other common divisor of $a$ and $b$, then $d'|d$. We write $d=gcd(a,b)$.
    \item We say that two integers $a$ and $b$ are relatively prime if $gcd(a,b)=1$.
  \end{itemize}

\end{definition}
Proof of the $gcd$ theorem uses the division algorithm:
\begin{align*}
  a &= q_1b+r_1\;, 0\leq r_1 < b\\
  b&=q_2r_1+r_2\;, 0\leq r_2<r_1, r_1\neq0\\
  r_1&=q_3r_2+r_3\;, 0\leq r_3<r_2, r_2\neq0\\
  &\vdots\\
  r_{n-2}&=q_nr_{n-1}+r_n\;, 0\leq r_n<r_{n-1}, r_{n-1}\neq0\\
  r_{n-1}&=q_{n+1}r_n+0.
\end{align*}
The first non-zero remainder is $r_n|gcd(a,b)$, or a common divisor. To show that it's the greatest common divisor, suppose $d=as'+bt'$ for some $s',t'\in\Z$ where $d=gcd(a,b)$. Then if $a=q_1b+r_1$ and $b=q_2r_1+0$, if $d>r$, but $d$ is a common divisor, then $a-q_1b=r_1$. Thus $d|a,d_b\implies d|r_1$, therefor $d$ can't be bigger than $r_1.$
\begin{example}
  Let $a=23, b=3$. Since $23=1\cdot3+20$, we have $q=1$ and $r=20$ which is not valid since $r$ is not less than $b$. We also have $23=8\cdot3-1$, so $q=8$ and $r=-1$ is not valid since $r$ is not nonnegative.

  The only valid one is $23=3\cdot7+2$, so $q=7$ and $r=2$. There is no other way to write $23=3\cdot q+r$ where $0\leq r<3$.
\end{example}

\begin{example}
  Use the Division Algorithm to find $gcd(2520, 378)$.
  \begin{align*}
    2520&=378(6)+252\\
    378&=252(1)+126\\
    252&=126(2)+0
  \end{align*}
  The last non-zero remainder is the gcd, so $gcd(2520, 378)=126$.
\end{example}
\section{Lecture 5 -- 1/28}
\begin{theorem}
  Let $p$ be an integer such that $p>1$. We say that $p$ is a \textbf{prime number}, or simply $p$ is \textbf{prime}, if the only positive divisors of $p$ are 1 and $p$ itself. An integer greater than 1 that is not prime is called a \textbf{composite number}.
\end{theorem}
Lemma. Euclid. Let $a$ and $b$ be integers and $p$ be a prime number. If $p|ab$, then either $p|a$ or $p|b$. Suppose $p|ab$ and suppose $p\nmid a$. Then $gcd(a,p)=1$. Thus there exist integers $s,t$ such that $as+pt=1$. Multiplying both sides by $b$ gives $abs+ptb=b$. Since $p|ab$ and $p|ptb$, we have $p|b$.
\begin{theorem}
  There exist an infinite number of primes.

  Assume $p_1, p_2, \cdots, p_n$ are all the primes. Let $N=p_1p_2\cdots p_n+1$. Then $N$ is either prime or composite. If $N$ is prime, then there exists a prime not in our list. If $N$ is composite, then it has a prime divisor $p$. Since $p$ divides $N$ and $p$ divides $p_1p_2\cdots p_n$, $p$ must also divide their difference, which is 1. This is a contradiction since no prime divides 1. Thus there exists a prime not in our list.
\end{theorem}
\subsection{3.1 Integer Equivalence Classes and Symmetries}
$\Z_n=\{[0], [1], [2],\cdots,[n-1]\}$ where arithmetic is mod $n$ and the equivalence classes are $a\equiv b(mod\ n)$ if $n|a-b$.
\begin{example}
  $\Z_6:$\\
  $[0]=\{0, \pm6, \pm12, \pm18\cdots\}=[6]=[-12]$ all have a remainder of 0 when divided by 6.\\
  $[1]=\{1, \pm5, \pm11, \pm17\cdots\}=[7]=[-11]$ all have a remainder of 1 when divided by 6.\\
  $\vdots$
\end{example}
\section{Lecture 6 -- 1/29}
\begin{definition}
  Let $\Z_n$ be the set of equivalence classes of the integers $mod\ n$ and $a,b,c\in\Z_n$.

  \begin{enumerate}
    \item Addition and multiplication are commutative. $ab\equiv ba$ and $a+b\equiv b+a$.
    \item Addition and multiplication are associate. $(ab)c\equiv a(bc)$ and $(a+b)+c\equiv a+(b+c)$.
    \item There are both additive and multiplicative identities. $0\in \Z_n$ has the property $a+0\equiv a$ and $1\in \Z_n$ has the property $a\cdot1\equiv a$.
    \item Multiplication distributes over addition. $a(b+c)\equiv ab+ac$.
    \item For every integer $a$ there is an additive inverse $-a$. $a+(n-a)\equiv (n-a)+a\equiv n\equiv0$.
    \item Let $a$ be a nonzero integer. Then $gcd(a,n)=1$ if and only if there exists a multiplicative $b$ inverse for $a(mod\ n)$.
  \end{enumerate}
\end{definition}
\begin{definition}
  A \textbf{symmetry} of a geometric figure is a rearrangement of the figure preserving the arrangement of its sides and vertices as well as its distances and angles. A map from the plane to itself preserving the symmetry of an object is called a \textbf{rigid motion}.
\end{definition}
\subsection{3.2 Definitions and Examples}
\begin{definition}
  A \textbf{binary operation} or \textbf{law of composition} on a set $G$ is a function $G\times G\to G$ that assigns to each pair $(a,b)\in G\times G$ a unique element $a \circ b$, or $ab$ in $G$, called the composition of $a$ and $b$.
\end{definition}
\begin{definition}
  A \textbf{group} is a set $G$ together with a law of composition $\circ$ that satisfies the following axioms:
  \begin{itemize}
    \item The law of composition is \textbf{associative}; that is, for all $a,b,c\in G$, we have $(a\circ b)\circ c=a\circ(b\circ c)$.
    \item There exists an \textbf{identity element} $e$ in $G$ such that for all $a\in G$, we have $e\circ a=a\circ e=a$.
    \item For each element $a\in G$, there exists an \textbf{inverse element} $a^{-1}\in G$ such that $a\circ a^{-1}=a^{-1}\circ a=e$.
  \end{itemize}
\end{definition}
A group with the property that $a\circ b=b\circ a$ for all $a,b\in G$ is called an \textbf{abelian group} or \textbf{commutative group}. Groups not satisfying this property are called \textbf{nonabelian} or \textbf{non-commutative}.

\section{Lecture 7 -- 2/3}

The identity element in a group $G$ is unique; that is, there exists only one element $e\in G$ such that $eg=ge=g$ for all $g\in G$.

If $g$ is any element in a group $G$, then the inverse of $g$, denoted by $g^{-1}$, is unique.

Let $G$ be a group. If $a,b\in G$, then $(ab)^{-1}=b^{-1}a^{-1}$.

Let $G$ be a group. For any $a\in G$, $(a^{-1})^{-1}=a$.

Let $G$ be a group and $a$ and $b$ be any two elements in $G$. Then the equations $ax=b$ and $xa=b$ have unique solutions in $G$.

If $G$ is a group and $a,b,c\in G$, then $ba=ca$ implies $b=c$ and $ab=ac$ implies $b=c$.

\begin{theorem}
  In a group, the usual law of exponents hold; that is, for all $g,h\in G$,
  \begin{enumerate}
    \item $g^mg^n=g^{m+n}$ for all $m,n\in\Z$.
    \item $(g^m)^n=g^{mn}$ for all $m,n\in\Z$.
    \item $(gh)^n=(h^{-1}g^{-1})^{-n}$ for all $n\in\Z$. Furthermore, if $G$ is abelian, then $(gh)^n=g^nh^n$.
  \end{enumerate}
\end{theorem}

\section{Lecture 8 -- 2/5}
\subsection{3.3 Subgroups}
We define a \textbf{subgroup} $H$ of a group $G$ to be a subset $H$ of $G$ such that when the group operation of $G$ is restricted to $H$, $H$ is a group in its own right.

Observe that every group with at least two elements will always have at least two subgroups, the subgroup consisting of the identity element alone and the entire group itself. The subgroup $H = \{e\}$ of a group is called the \textbf{trivial subgroup}. A subgroup that is a proper subset of is called a \textbf{proper subgroup}.
\begin{example}
  Let $G=(\Z,+)$, and $H=(\Z+).$ $H$ is a trivial subgroup of $G$. Then let $H'=(\{0\},+).$ $H'$ is the trival subgroup of $G$.
\end{example}
Since $\Z\subset\Q\subset\R\subset\C$, we have: $(\C,+)$ is a group, $(\R,+)$ is a subgroup of $(\C,+)$. If $H=(\R,+)$ is a subgroup of $G=(\C,+)$, then consider $(a_1+b_1i) + (a_2+b_2i) \in G \implies (a_1+a_2)+(b_1+b_2)i\in G.$

Let $H$ be a subset of group $G$. Then $H$ is a subgroup of $G$ if and only if
\begin{itemize}
  \item \textbf{Closure}: For all $a,b\in H$, $ab\in H$.
  \item \textbf{Identity}: The identity element $e$ of $G$ is in $H$.
  \item \textbf{Inverses}: For all $a\in H$, $a^{-1}\in H$.
\end{itemize}
\subsection{4.1 Cyclic Subgroups}
\begin{theorem}
  Let $G$ be a group and $a$ be any element in $G$. Then the set $\langle a \rangle = \{a^n:n\in \Z\}$ is a subgroup of $G$. Furthermore, $\langle a \rangle$ is the smallest subgroup of $G$ that contains $a$ and is called the \textbf{cyclic subgroup} of $G$ generated by $a$.
\end{theorem}
\section{Lecture 9 -- 2/10}
For $a\in G$, we call $\langle a \rangle$ the \textbf{cyclic subgroup} of $G$ generated by $a$.

If $G$ contains some element $a$ such that $G=\langle a \rangle$, then $G$ is a \textbf{cyclic group} and $a$ is called a \textbf{generator} of $G$.

If $a$ is an element of a group $G$, we define the \textbf{order} of $a$ to be the smallest positive integer $n$ such that $a^n=e$, and we write $|a|=n$. If there is no such integer $n$, we say that the order of $a$ is infinite and write $|a|=\infty$ to denote the order of $a$.
\begin{theorem}
  Every cyclic group is abelian.

  Also, every subgroup of a cyclic group is cyclic.

  Corollary: The subgroups of $\Z$ are exactly $n\Z$ for $n=0,1,2,\cdots$.
\end{theorem}
\section{Lecture 10 -- 2/11}
Proposition: Let $G$ be a cyclic group of order $n$ and suppose that $a$ is a generator for $G$. Then $a^k = e$ if and only if $n|k$, or $k$ divides $n$.
\begin{theorem}
  Let $G$ be a cyclic group of order $n$ and suppose that $a$ is a generator of the group. If $b=a^k$, then the order of $b$ is $\frac{n}{d}$, where $d=gcd(k,n)$.

  Corollary: The generators of $\Z_n$ are the integers $r$ such that $1\leq r<n$ and $gcd(r,n)=1$.
\end{theorem}
\section{Lecture 11 -- 2/12}
\subsection{4.2 Multiplicative Group of Complex Numbers}
$\C=\{x+oy:x,y\in\R,i^2=-1\}.$

If $z=x+iy$, then $z=r(\cos{\theta}+i\sin{\theta})=cis\theta=re^{i\theta}$ where $r=\sqrt{x^2+y^2}$ and $\theta=\tan^{-1}{\frac{y}{x}}$.

The circle group $\T=\{z\in\C:|z|=1\}$ is a subgroup of $\C^*$.
\begin{theorem}
  If $z^n=1$, then the $n^{\text{th}}$ roots of unity are $z=\cos\frac{2k\pi}{n}+i\sin\frac{2k\pi}{n}$ for $k=0,1,2,\cdots,n-1$. Furthermore, the $n^{\text{th}}$ roots of unity form a cyclic subgroup of $\T$ of order $n$.
\end{theorem}
\subsection{5.1 Definitions and Notation}
\begin{theorem}
  The symmetric group on $n$ letters, $S_n$, is a group with $n!$ elements, where the binary operation is the composition of maps.
\end{theorem}
\textbf{Cycle Notation:} A permutation $\sigma\in S_X$ is a cycle of length $k$ if there exist elements $a_1, a_2, \cdots, a_k\in X$ such that $\sigma(a_1)=a_2, \sigma(a_2)=a_3, \cdots, \sigma(a_k)=a_1$, and $\sigma(x)=x$ for all other $x\in X$. We write $(a_1a_2\cdots a_k)$ to denote the cycle $\sigma$.

A subgroup of $S_n$ is called a \textbf{permutation group}.

Two cycles in $S_n, \sigma=(a_1a_2\cdots a_k)$ and $\tau=(b_1b_2\cdots b_m)$ are \textbf{disjoint} if $a_i\neq b_j$ for all $i$ and $j$.

Let $\sigma$ and $\tau$ be two disjoint cycles in $S_n$. Then $\sigma\tau=\tau\sigma$.

Every permutation in $S_n$ can be written as the product of disjoint cycles.

\section{Lecture 12 -- 2/18}
\begin{itemize}
  \item A cycle of length 2 is called a \textbf{transposition}.
  \item \textbf{Proposition:} Any permutation of a finite set containing at least two elements can be written as the product of transpositions.
  \item If the identity $id$ is written as the product of $r$ transpositions, $id=\tau_1\tau_2\cdots\tau_r$, then $r$ is even.
\end{itemize}
\begin{theorem}
  If a permutation $\sigma$ can be expressed as the product of an even number of transpositions, then any other product of transpositions equaling must also contain an even number of transpositions.

  We define a permutation to be even if it can be expressed as an even number of transpositions and odd if it can be expressed as an odd number of transpositions.
\end{theorem}
\begin{theorem}
  The set of even permutations $A_n$ is a subgroup of $S_n$.
\end{theorem}
\textbf{Proposition:} The number of even permutations is $S_n, n\geq2$, is equal to the number of odd permutations; hence the order of $A_n$ is $\frac{n!}{2}.$
\section{Lecture 13 -- 2/19}
\subsection{5.2 Dihedral Groups}
For $n=3,4,\cdots,$ we define the \textbf{nth dihedral group} to be the group of rigid motions of a regular $n$-gon.
\begin{theorem}
  The dihedral group, $D_n$, is a subgroup of $S_n$ of order $2n$.
\end{theorem}
\begin{theorem}
  The group $D_n, n\geq3$, consists of all products of the elements $r$ and $s$, where $r$ has order $n$ and $s$ has order 2, and these two elements satisfy the relation $srs=r^{-1}$.
\end{theorem}
\subsection{6.1 Cosets}
Let $G$ be a group and $H$ be a subgroup of $G$. Define a \textbf{left coset} of with representative $g\in G$ to be the set $gH=\{gh:h\in H\}$. Define a \textbf{right coset} of $H$ with representative $g$ to be the set $Hg=\{hg:h\in H\}$.
\section{Lecture 14 -- 2/26}
\textbf{Lemma: } Let $H$ be a subgroup of a group $G$ and suppose that $g_1,g_2\in G$. The followng conditions are equivalent:
\begin{itemize}
  \item $g_1H=g_2H$
  \item $Hg_1^{-1}=Hg_2^{-1}$
  \item $g_1H\subset g_2H$
  \item $g_2\in g_1H$
  \item $g_1^{-1}g_2\in H$.
\end{itemize}
\begin{theorem}
  Let $H$ be a subgroup of a group $G$. Then the left cosets of $H$ in $G$ form a partition of $G$. That is, the group $G$ is the disjoint union of the left cosets of $H$ in $G$.
\end{theorem}
\subsection{6.2 Lagrange's Theorem}
\textbf{Proposition: }Let $H$ be a subgroup of $G$ with $g\in G$ and define a map $\phi: H\to gH$ by $\phi(h)=gh$. The map is bijective; hence, the number of elements in $H$ is the same as the number of elements in $gH$.

\begin{theorem}
  Let $G$ be a finite group and let $H$ be a subgroup of $G$. Then $|G|/|H|=[G:H]$ is the number of distinct left cosets of $H$ in $G$, called the \textbf{index} of $H$ in $G$. In particular, the number of elements in $H$ must divide the number of elements in $G$.
\end{theorem}

\textbf{Corollary:} Suppose that $G$ is a finite group and $g\in G$. Then the order of $g$ must divide the number of elements in $G$.

\textbf{Corollary:} Let $|G|=p$ with $p$ a prime number. Then $G$ is cyclic and any $g\in G$ such that $g\neq e$ is a generator.

\textbf{Corollary:} Let $H$ and $K$ be subgroups of a finite group $G$ such that $K\subset H \subset G$. Then $[G:K]=[G:H][H:K]$.
\subsection{6.3 Fermat's and Euler's Theorems}
The \textbf{Euler $\phi$-function} is the map $\phi:\N\to\N$ defined by $\phi(1)=1$, and, for $n>1$, $\phi(n)$ is the number of positive integers $m$ with $1\leq m<n$ and $gcd(m,n)=1$.

\section{Lecture 15 -- 3/3}
\begin{theorem}
  Let $U(n)$ be the group of units in $\Z_n$ (units are elements with multiplicative inverses, aka coprime to $n$) under multiplication. Then $|U(n)|=\phi(n)$.
\end{theorem}
\begin{theorem}
  Euler's theorem: Let $a$ and $n$ be integers such that $n>0$ and $gcd(a,n)=1$. Then $a^{\phi(n)}\equiv 1\mod n$.
\end{theorem}
\begin{example}
  What is the last digit of $457^{327}$?

  Since $457\equiv 7 \mod 10\implies 457^{327}\equiv y^{327}\mod 10$. Since $gcd(7,10)=1$, we can apply Euler's theorem to get $7^{\phi(10)}\equiv 1\mod 10$. Since $\phi(10)=4$, we have $7^4\equiv 1\mod 10$. Thus $7^{327}\equiv 7^{4\cdot81+3}\equiv (7^4)^{81}\cdot7^3\equiv 1^{81}\cdot343\equiv 3\mod 10$. Therefore the last digit of $457^{327}$ is 3.
\end{example}
\begin{theorem}
  Fermat's Little Theorem: Let $p$ be any prime number and suppose that $p\nmid a$($p$ does not divide $a$). Then $a^{p-1}\equiv 1\mod p$. Furthermore, for any integer $b$, $b^p\equiv b\mod p$.
\end{theorem}
\subsection{9.1 Isomorphisms: Definition and Examples}
Two groups $(G,\cdot)$ and $(H,\circ)$ are \textbf{isomorphic} if there exists a one-to-one and onto map $\phi:G\to H$ such that the group operation is preserved; that is, $\phi(a\cdot b)=\phi(a)\circ\phi(b)$ for all $a,b\in G$. If $G$ is isomorphic to $H$, we write $G\cong H$. The map is called an isomorphism.
\section{Lecture 16 -- 3/4}
\begin{theorem}
  All cyclic groups of infinite order are isomorphic to $\Z$.
\end{theorem}
\begin{theorem}
  If $G$ is a cyclic group of order $n$, then $G$ is isomorphic to $\Z_n$.

  Corollary: If $G$ is a group of order $p$, where $p$ is a prime number, then $G$ is isomorphic to $\Z_p$.
\end{theorem}
\begin{theorem}
  Cayley: Every group is isomorphic to a group of permutations.
\end{theorem}
\subsection{9.2 Direct Products}
If $(G,\cdot)$ and $(H,\circ)$ are groups, then we can make the Cartesian product of $G$ and $H$ into a new group. As a set, our group is just the ordered pairs $(g,h)\in G\times H$. We can define a binary operation on $G\times H$ by $(g_1,h_1)(g_2,h_2)=(g_1\cdot g_2, h_1\circ h_2)$.

\textbf{Proposition:} Let $G$ and $H$ be groups. The set $G\times H$ is a group. The group $G\times H$ is called the \textbf{external direct product} of $G$ and $H$.

\begin{theorem}
  Let $(g,h)\in G\times H$. If $g$ and $h$ have finite orders $r$ and $s$ respectively, then the order of $(g,h)$ in $G\times H$ is the least common multiple of $r$ and $s$.
\end{theorem}
\begin{theorem}
  The group $\Z_m\times\Z_n$ is isomorphic to $\Z_{mn}$ if and only if $gcd(m,n)=1$.
\end{theorem}
\section{Lecture 17 -- 3/5}
\textbf{Corollary:} Let $n_1, \dots,n_k$ be positive integers. Then
\[\Pi_{i=1}^k \Z_{n_i}\cong \Z_{n_1\cdots n_k}\]
if and only if $gcd(n_i,n_j)=1$ for $i\neq j$.

\textbf{Corollary:} If $m={p_1}^{e_1}\dots p_k^{e_k}$ where the $p_i$s are distinct primes, then
\[\Z_m\cong\Z_{p_1^{e_1}}\times\cdots\times\Z_{p_k^{e_k}}.\]

\begin{definition}
  Let $G$ be a group with subgroups $H$ and $K$ satisfying the following conditions:
  \begin{itemize}
    \item $G=HK=\{hk:h\in H,j\in K\}$,
    \item $H\cap K=\{e\}$,
    \item $hk=kh$ for all $h\in H$ and $k\in K$.
  \end{itemize}
  Then $G$ is the \textbf{internal direct product} of $H$ and $K$.
\end{definition}

\begin{theorem}
  Let $G$ be the internal direct product of subgroups $H$ and $K$. Then $G$ is isomorphic to $H\times K$.
\end{theorem}
\section{Lecture 17 -- 3/17}
\subsection{10.1 Factor Groups and Normal Subgroups}
A subgroup $H$ of a group $G$ is \textbf{normal} in $G$ if $gH=Hg$ for all $g\in G$.
\begin{example}
  Consider $3\Z$ in $\Z$.

  $H=3\Z=\{\cdots,-6,--3,0,3,6,\cdots\}=\{3n:n\in\Z\}.$

  $0+H=3+H=6+H=\cdots=H+0=H+3=\cdots$
  $1+H=4+H=7+H=\cdots=H+1=H+4=\cdots$
  $2+H=5+H=8+H=\cdots=H+2=H+5=\cdots$

  If $G$ is abelian, then every subgroup is normal.

\end{example}

\begin{example}
  Consider $H = \{e, (12)\}$ in $S_3$.

  Note $eH=He=H$ and $(12)H=\{(12), e\}=\{e, (12)\}=H(12)$.

  Consider $(123)H=\{(123)e, (123)(12)\}=\{(123), (13)\}$ and $H(123)=\{e(123), (12)(123)\}=\{(123), (23)\}$.

  $H(123)=\{(123), (23)\}\neq\{(123), (13)\}=(123)H$. Thus $H$ is not normal in $S_3$.
\end{example}

\begin{example}
  Consider $A_3$ in $S_3$.

  $A_3=\{e, (123), (132)\}$. We know that $eA_3=A_3e=A_3$. Also, $(123)A_3=\{(123), (132), e\}=A_3(123)$ and $(132)A_3=\{(132), e, (123)\}=A_3(132)$.

  Consider $(12)A_3=\{(12), (13), (23)\}$ and $A_3(12)=\{(12), (13), (23)\}$.

  So $(12)A_3=A_3(12)$, and $A_3$ is normal in $S_3$.
\end{example}

\begin{theorem}
  Let $G$ be a group and $N$ be a subgroup of $G$. Then the following statements are equivalent.
  \begin{itemize}
    \item The subgroup $N$ is normal in $G$.
    \item For all $g\in G$, $gNg^{-1}=N$.
    \item For all $g\in G$, $gNg^{-1}\subset N$.
  \end{itemize}
\end{theorem}

\begin{definition}
  If $N$ is a normal subgroup of a group $G$, then the cosets of $N$ in $G$ form a group $G/N$ under the operation $(aN)(bN)=abN.$ This group is called the \textbf{factor} or \textbf{quotient group} of $G$ and $N$.
\end{definition}
\begin{example}
  What is the factor group $\Z/3\Z$?

  $\Z/3\Z=\{0+3\Z,1_3\Z,2+3\Z\}=\\\{\{\dots,-6,-3,0,3,6,\dots\}, \{...,-5,-2,1,4,7,...\}, \{...,-4,-1,2,5,8,...\}\}$.

  $(0+3\Z)+(1+3\Z)=(0+1)+3\Z=1+3\Z$.

  $(1+3\Z)+(2+3\Z)=(1+2)+3\Z=0+3\Z$.

  $(2+3\Z)+(2+3\Z)=(2+2)+3\Z=1+3\Z$.
\end{example}

\begin{theorem}
  Let $N$ be a normal subgroup of a group $G$. The cosets of $N$ in $G$ form a group $G/N$ of order $[G:N]$.
\end{theorem}
\section{Lecture 18 -- 3/18}
\subsection{11.1 Group Homomorphisms}
A \textbf{homomorphism} between groups $(G,\cdot)$ and $(H,\circ)$ is a map $\phi:G\to H$ such that $\phi(g_1\cdot g_2)=\phi(g_1)\circ\phi(g_2)$ for $g_1,g_2\in G$. The range of $\phi$ is  called the \textbf{homomorphic image} of $G$.
\begin{example}
  Let $\phi:\Z\to\Z_4\times\Z_9$ be a homomorphism with $\phi(1)=(2,3).$ Find $\phi(n)$ for any $n\in\Z$. Which elements $n\in\Z$ have $\phi(n)=(0,0)$?

  $\phi(2)=\phi(1+1)=\phi(1)+\phi(1)=(2,3)+(2,3)=(0,6)\in\Z_4\times\Z_9$.

  $\phi(3)=\phi(1)+\phi(1)+\phi(1)=\phi(2)+\phi(1)=(0,6)+(2,3)=(2,0)$.

  $\phi(4)=(0,3)$, $\phi(5)=(2,6)$, $\phi(6)=(0,0)$.
\end{example}

Let $\phi:G_1\to G_2$ be a homomorphism of groups. Then
\begin{enumerate}
  \item If $e$ is the identity of $G_1$, then $\phi(e)$ is the identity of $G_2$.
  \item For any element $g\in G,\phi(g^{-1})=[\phi(g)]^{-1}$.
  \item If $H_1$ is a subgroup of $G_1$, then $\phi(H_1)$ is a subgroup of $G_2$.
  \item If $H_2$ is a subgroup of $G_2$, then $\phi^{-1}(H_2)=\{g\in G_1:\phi(g)\in H_2\}$ is a subgroup of $G_1$; Furthermore, if $H_2$ is normal in $G_2$, then $\phi^{-1}(H_2)$ is normal in $G_1$.
\end{enumerate}

Let $\phi:G\to H$ be a group homomorphism and suppose that $e$ is the identity of $H$. The subgroup $\phi^{-1}(\{e\})$ is called the \textbf{kernel} of $\phi$ and is denoted by $ker(\phi)$.
\begin{theorem}
  Let $\phi:G\to H$ be a group homomorphism. Then the kernel of $\phi$ is a normal subgroup of $G$.
\end{theorem}
\section{Lecture 19 -- 3/19}
\subsection{11.2 The Isomorphism Theorems}
Let $H$ be a normal subgroup of $G$. Define the \textbf{natural} or \textbf{canonical homomorphism} $\phi:G\to G/H$ by $\phi(g)=gH$.
\begin{example}
  Let $G=\Z$ and $H=5\Z$.

  Let $\phi:\Z\to\Z/5\Z$, $\phi(n)=n+5\Z$.

  $\phi(0)=0+5\Z=[0]$, $\phi(1)=1+5\Z=[1]$, $\phi(2)=2+5\Z=[2]$, $\phi(3)=3+5\Z=[3]$, $\phi(4)=4+5\Z=[4]$, $\phi(5)=0+5\Z=[0]$. Note: $ker(\phi)=5\Z$ (kernel: a group and happens to be the image of $\phi$).
\end{example}
\begin{theorem}
  \textbf{First Isomorphism Theorem:} If $\psi:G\to H$ is a group homomorphism with $K=ker\psi$, then $K$ is normal in $G$. Let $\phi:G\to G/K$ be the canonical homomorphism. Then there exists a unique isomorphism $\eta:G/K\to\psi(G)$ such that $\psi=\eta\phi$.
\end{theorem}
\begin{example}
  Let $\psi:\Z\to\Z_6$ be a group homomorphism with $\psi(1)=2$. Verify the First Isomorphism Theorem.

  $\psi(n)=2n\mod 6.$ $K=ker\psi=3\Z\subset\Z$.

  Let $\phi:\Z\to\Z/ker\psi\implies\phi:\Z\to\Z/3\Z$ be the canonical homomorphism. Then the First Iso. Theorem $\implies \Z/3\Z\cong$ image of $\psi$.
\end{example}
\begin{theorem}
  \textbf{Second Isomorphism Theorem:} Let $H$ be a subgroup of a group $G$ (not necessarily normal in $G$) and $N$ a normal subgroup of $G$. Then $HN$ is a subgroup of $G, H\cap N$ is a normal subgroup of $H$, and $HN/N\cong H/H\cap N$.
\end{theorem}
\begin{example}
  Let $H=6\Z$ and $N=10\Z$ be subgroups of $G=\Z$. Verify the Second Isomorphism Theorem.

  $HN=\{h+n:h\in H, n\in N\}\implies(6\Z)(10\Z)=\{(6m)+(10n):m,n\in\Z\}=2\Z.$

  $H\cap N=\{m:m=6k,m=10n\text{ for some }k,n\in\Z\}=30\Z.$

  $H/H\cap N=6\Z/30\Z=\{[0],[6], [12], [18], [24]\}$.

  $HN/N=2\Z/10\Z=\{[0],[2],[4],[6],[8]\}$.
\end{example}
\section{Lecture 20 -- 3/24}
\begin{theorem}
  \textbf{Correspondence Theorem:} Let $N$ be a normal subgroup of a group $G$. Then $H\mapsto H/N$ is a one-to-one correspondence between the set of subgroups $H$ of $G$ containing $N$ and the set of subgroups of $G/N$. Furthermore, the normal subgroups of $G$ containing $N$ correspond to normal subgroups of $G/N$.
\end{theorem}
\begin{theorem}
  \textbf{Third Isomorphism Theorem:} Let $G$ be a group and $N$ and $H$ be normal subgroups of $G$ with $N\subset H$. Then $G/H\cong\frac{G/N}{H/N}.$
\end{theorem}
\begin{example}
  Let $G=\Z,H=4\Z,$ and $N=24\Z$. Verify the Third Isomorphism Theorem.

  $24\Z=N\subset 4\Z=H\subset\Z=G$.

  $G/H=\Z/4\Z\cong\Z_4$.

  $G/N=\Z/24\Z\cong\Z_{24}$.

  $H/N=4\Z/24\Z\cong\Z_6$.

  $\frac{G/N}{H/N}=\frac{\Z/24\Z}{4\Z/24\Z}\cong\Z_4$.
\end{example}
\section{Lecture 21 -- 3/25}
\subsection{13.1 Finite Abelian Groups}
\begin{itemize}
  \item Suppose that $G$ is a group and let $\{g_i\}$ be a set of elements in $G$, where $i$ is in some index set $I$ (not necessarily finite). The smallest subgroup of containing all of the $g_i$'s is the subgroup of $G$ \textbf{generated} by the $g_i$'s.
  \item If this subgroup of $G$ is in fact all of $G$, then $G$ is \textbf{generated} by the set $\{g_i:i\in I\}$. In this case the $g_i$'s are called \textbf{generators} of $G$.
  \item If there is a finite set $\{g_i:i\in I\}$ that generates $G$, then $G$ is \textbf{finitely generated}.
  \item Define a group $G$ to be a $p$-group if every element in $G$ has as its order a power of $p$, where $p$ is a prime number.
\end{itemize}
\begin{theorem}
  \textbf{Fundamental Theorem of Finite Abelian Groups:} Every finite abelian group $G$ is isomorphic to a direct product of cyclic groups of the form \[\Z_{p_1^{\alpha_1}}\times\Z_{p_2^{\alpha_2}}\times\cdots\times\Z_{p_n^{\alpha_n}}\]where the $p_i$'s are primes (not necessarily distinct).
\end{theorem}
\begin{theorem}
  \textbf{Fundamental Theorem of Finitely Generated Abelian Groups:} Every finitely generated abelian group is isomorphic to a direct product of cyclic groups of the form \[\Z_{p_1^{\alpha_1}}\times\Z_{p_2^{\alpha_2}}\times\cdots\times\Z_{p_n^{\alpha_n}}\times\Z\cdots\times\Z\] where the $p_i$'s are primes (not necessarily distinct).
\end{theorem}

Many of the sets we have explored have 2 operations: addition and multiplication. For example, $\Z,\Q,\R,\C,\Z_n,M_2(\R)=\{2\times2\text{ matrices in }\R\}$, etc. Under addition, all of these sets are abelian, and most are groups. Under multiplication, most of these sets are not groups.

\subsection{16.1 Rings}
A nonempty set $R$ is a ring if it has two closed binary operations, addition and multiplication, satisfying the following conditions:
\begin{itemize}
  \item $a+b=b+a$ for $a,b\in R$.
  \item $a+(b+c)=(a+b)+c$ for $a,b,c\in R$.
  \item There is an element $0\in R$ such that $a+0=0+a$ for all $a\in R$.
  \item For every element $a\in R$, there is an element $-a\in R$ such that $a+(-a)=0$.
  \item $(ab)c=a(bc)$ for all $a,b,c\in R$.
  \item $a(b+c)=ab+ac$ and $(a+b)c=ac+bc$ for $a,b,c\in R$.
\end{itemize}

If there is an element $1\in R$ such that $1\neq0$ and $1a=a1=a$ for each element $a\in R$, we say that $R$ is a ring with \textbf{unity} or \textbf{identity}.

A ring $R$ for which $ab=ba$ for all $a,b\in R$ is called a \textbf{commutative ring}.

A commutative ring with identity is called an \textbf{integral domain} if, for ever $a,b\in R$ such that $ab=0$, either $a=0$ or $b=0$.

A \textbf{division ring} is a ring $R$, with an identity, in which every nonzero element is a unit; that is, for each $a\in R$ with $a\neq0$, there exists a unique element $a^{-1}\in R$ such that $aa^{-1}=a^{-1}a=1$.

A commutative division ring is called a \textbf{field}.

\section{Lecture 22 -- 3/31}
\subsection{16.2 Integral Domains and Fields}
\textbf{Cancellation Law:} Let $D$ be a commutative ring with identity. Then $D$ is an integral domain if and only if for all nonzero elements $a\in D$ with $ab=ac$, we have $b=c$.
\begin{theorem}
  Every finite integral domain is a field.

  Let $D=\{0,1,a_1,a_2,\cdots,a_n\}$ be a finite
\end{theorem}

For any nonnegative integer $n$ and any element $r$ in a ring $R$ we write $r+r+\cdots+r$ ($n$ times) as $nr$. We define the \textbf{characteristic} of a ring $R$ to be the least positive integer $n$ such that $nr=0$ for all $r\in R$. If no such integer exists, then the characteristic of $R$ is defined to be $0$. We will denote the characteristic of $R$ by $char(R)$.

\textbf{Lemma:} Let $R$ be a ring with identity. If 1 has order $n$, then the characteristic of $R$ is $n$.
\begin{theorem}
  The characteristic of an integral domain is either prime or zero.
\end{theorem}

Every ring is an abelian group under addition.

If $R$ is a finite ring, then we know the structure under addition, but there may be several types of multiplication.

\subsection{16.3 Ring Homomorphisms and Ideals}
An \textbf{ideal} in a ring $R$ is a subring $I$ of $R$ such that if $a$ is in $I$ and $r$ is in $R$, then both $ar$ and $ra$ are in $I$; that is $rI\subset I$ and $Ir\subset I$ for all $r\in R$.

If $R$ and $S$ are rings, then a \textbf{ring homomorphism} is a map $\phi:R\to S$ satisfying \[\phi(a+b)=\phi(a)+\phi(b)\]\[\phi(ab)=\phi(a)\phi(b)\] for all $a,b\in R$.

If $\phi$ is a one-to-one and onto homomorphism, then $\phi$ is called an \textbf{isomorphism} of rings.

\section{Lecture 23 -- 4/1}
\begin{theorem}
  \textbf{First Isomorphism Theorem for Rings:} Let $\psi:R\to S$ be a ring homomorphism. Then $ker\psi$ is an ideal of $R$. If $\phi:R\to R/ker\psi$ is the canonical homomorphism, then there exists a unique isomorphism $\eta:R/ker\psi\to im\psi$ such that $\psi=\eta\phi$.
\end{theorem}
\section{Lecture 24 -- 4/2}
\begin{theorem}
  \textbf{Second Isomorphism Theorem for Rings:} Let $I$ be a subring of a ring $R$ and let $J$ be an ideal of $R$. Then $I\cap J$ is an ideal of $I$ and $I/I\cap J\cong (I+J)/J$.
\end{theorem}
\begin{theorem}
  \textbf{Third Isomorphism Theorem for Rings:} Let $R$ be a ring and  $I$ and $J$ be ideals of $R$ where $J\subset I$. Then $R/J\cong\frac{R/J}{I/J}$.
\end{theorem}
\begin{theorem}
  \textbf{Correspondence Theorem for Rings:} Let $I$ be an idea of a ring $R$. Then $S\mapsto S/I$ is a one-to-one correspondence between the set of subrings $S$ containing $I$ and the set of subrings of $R/I$. Futhermore, the ideals of $R$ containing $I$ correspond to ideals of $R/I.$
\end{theorem}
A proper ideal $M$ of a ring $R$ is a \textbf{maximal ideal} of $R$ if the ideal $M$ is not a proper subset of any ideal of $R$ except $R$ itself. That is, $M$ is a maximal ideal if for any ideal $I$ properly containing $M$, $I = R$.
\begin{theorem}
  Let $R$ be a commutative ring with identity and $M$ an ideal in $R$. Then $M$ is a maximal ideal of $R$ is and only if $R/M$ is a field.
\end{theorem}
\subsection{17.1 Polynomial Rings}
\begin{definition}
  Throughout this chapter, $R$ denotes a dommutative ring with identity.
  \begin{itemize}
    \item An expression of the form \[f(x)=\sum_{i=1}^n a_ix^i=a_0+a_1x+a_2x^2+\cdots+a_nx^n\] where $a_i\in R$ and $a_n\neq0$, is called a \textbf{polynomial over $R$} with \textbf{indeterminate $x$}.
    \item The elements $a_0, a_1,\cdots,a_n$ are called the \textbf{coefficients} of $f$, and $a_n$ is called the \textbf{leading coefficient}.
    \item A polynomial is \textbf{monic} if its leading coefficient is 1.
    \item If $n$ is the largest nonnegative integer such that $a_n\neq0$, we say the \textbf{degree} of $f$ is $n$ and write $deg\;f(x)=n$.
    \item If no such $n$ exists (i.e, $f=0$ is the zero polynomial), the degree of $f$ is defined to be $-\infty$.
    \item The set of all polynomials with coefficients in $R$ is denoted $R[x]$.
    \item Two polynomials $p(x)=a_0+a_1x+\cdots+a_nx^n$ and $q(x)=b_0+b_1x+\cdots+b_mx^m$ are \textbf{equal} if $a_i=b_i$ for all $i\geq0$.
  \end{itemize}
\end{definition}

Let $p(x)=a_0+a_1x+\cdots+a_nx^n$ and $q(x)=b_0+b_1x+\cdots+b_mx^m$ be polynomials in $R[x]$. We define the sum of $p$ and $q$ to be the polynomial \[p(x)+q(x)=(a_0+b_0)+(a_1+b_1)x+\cdots+(a_k+b_k)x^k\] where $k=max(n,m)$ and we let $a_i=0$ for $i>n$ and $b_j=0$ for $j>m$. We define the product of $p$ and $q$ to be the polynomial \[p(x)q(x)=\sum_{i=0}^{n+m}c_ix^i\] where $c_i=\sum_{j=0}^ia_jb_{i-j}$.

\begin{theorem}
  Let $R$ be a commutative ring with identity. Then $R[x]$ is a commutative ring with identity.
\end{theorem}
\begin{theorem}
  Let $p(x)$ and $q(x)$ be polynomials in $R[x]$, where $R$ is an integral domian. Then \[deg(p(x)q(x))=deg\;p(x)+deg\;q(x).\] Furthermore, $R[x]$ is an integral domain.
\end{theorem}

\section{Lecture 25 -- 4/7}
\subsection{17.2 The Division Algorithm}
\begin{theorem}
  Let $f(x)$ and $g(x)$ be poliynomials in $F[x]$, where $F$ is a field and $g(x)$ is a nonzero polynomial. Then there exist unique polynomials $q(x),r(x)\in F[x]$ such that $f(x)=g(x)q(x)+r(x)$, where either $deg\;r(x)<deg\;g(x)$ or $r(x)=0$.
\end{theorem}

\textbf{Zero/Root of a Polynomial:} Let $p(x)\in F[x]$, where $F$ is a field, and let $\phi_a:$

\begin{theorem}
  \textbf{Factor Theorem:} Let $F$ be a field. An element $\alpha\in F$ is a zero of $p(x)\in F[x]$ if and only if $(x-\alpha)$ is a factor of $p(x)$ in $F[x]$.
\end{theorem}

\begin{theorem}
  \textbf{Bound on Number of Zeros:} Let $F$ be a field. A non-zero polynomial $p(x)$ of degree $n$ in $F[x]$ has at most $n$ distinct zeros in $F$.
\end{theorem}

\textbf{Greatest Common Divisor of Polynomials:} Let $F$ be a field. A monic polynomial $d(x)\in F[x]$ is a \textbf{greatest common divisor} of $p(x),q(x)\in F[x]$ if $d(x)$ evenly divides both $p(x)$ and $q(x)$, and if for any other polynomial $d'(x)$ dividing both $p(x)$ and $q(x)$, we have $d'(x)|d(x)$.
We write $d(x)=gcd(p(x),q(x))$. Two polynomials $p(x)$ and $q(x)$ are \textbf{relatively prime} if $gcd(p(x),q(x))=1$.

\begin{theorem}
  \textbf{Existence and Uniqueness of the GCD: }Let $F$ be a field and supppose $d(x)$ is a greatest common divisor of $p(x)$ and $q(x)$ in $F[x]$. Then there exist polynomials $r(x)$ and $s(x)$ such that $d(x)=r(x)p(x)+s(x)q(x)$. Furthermore, the greatest common divisor or two polynomials is unique.
\end{theorem}

\subsection{17.3 Irreducible Polynomials}
A nonconstant polynomial $f(x)\in F[x]$ is \textbf{irreducible over a field} $F$ if $f(x)$ cannot be expressed as a product $g(x)h(x)$ where $g(x), h(x)\in F[x]$ and the degrees of $g(x)$ and $h(x)$ are both strictly less than $deg\;f(x)$. Irreducible polynomials function as the "prime numbers" of polynomial rings.

\begin{theorem}
  \textbf{Rational Coefficient Decomposition: }Let $p(x)\in\Q[x]$. Then \[p(x)=\frac{r}{s}(a_0+a_1x+\cdots+a_nx^n)\] where $r,s,a_9,\dots,a_n$ are integerts, the $a_i$'s are relative prime, and $gcd(r,s)=1$.
\end{theorem}

\begin{theorem}
  \textbf{Gauss's Lemma: }Let $p(x)\in\Z[x]$ be a monic polynomial such that $p(x)$ factors into a product of two polynomials $\alpha(x)$ and $\beta(x)$ in $\Q[x]$, where the degrees of both $\alpha(x)$ and $\beta(x)$ are less than the degree of $p(x)$. Then $p(x)=a(x)b(x)$ where $a(x)$ and $b(x)$ are monic polynomials in $\Z[x]$ with $deg\;\alpha(x)=deg\;a(x)$ and $deg\;\beta(x)=deg\;b(x)$.
\end{theorem}
\section{Lecture 26 -- 4/8}
A proper ideal $P$ in a commutative ring $R$ is called a \textbf{prime ideal} if whenever $ab\in P$, then either $a\in P$ or $b\in P$.

Let $R$ be a commutative ring with identity 1, where $1\neq0$. Then $P$ is a prime ideal in $R$ if and only if $R/P$ is an intcegral domain.

\textbf{Principal Ideal in $F[x]$:} Let $F$ be a field. Recall that a \textbf{principal ideal} in $F[x]$ is an ideal $\langle p(x)\rangle$ generated by some polynomial $p(x)$; that is, $\langle p(x)\rangle=\{p(x)q(x):q(x)\in F[x]\}$.

\begin{theorem}
  \textbf{Every Ideal in $F[x]$ is Principal:} If $F$ is a field, then every ideal in $F[x]$ is a principal ideal.
\end{theorem}
\begin{theorem}
  \textbf{Maximal Ideals and Irreducible Polynomials:} Let $F$ be a field and $p(x)\in F[x]$. Then $\langle p(x)\rangle$ is a maximal ideal of $F[x]$ if and only if $p(x)$ is irreducible over $F$.
\end{theorem}
\begin{theorem}
  \textbf{Integer Zeros of Monic Integer Polynomials:} Let $p(x)=x^n+a_{n-1}x^{n-1}+\cdots+a_0\in\Z[x]$ with $a_0\neq0$. If $p(x)$ has a zero in $\Q$, then it also has a zero $\alpha\in\Z$, and moreover $\alpha|a_0$.
\end{theorem}

\begin{theorem}
  \textbf{Eisenstein's Criterion:} Let $p$ be a prime and suppose $f(x)=a_nx^n+\cdots+a_0\in\Z[x]$. If $p|a_i$ for $i=0,1,\dots,n-1$, but $p\nmid a_n$ and $p^2\nmid a_0$, then $f(x)$ is irreducible over $\Q$.
\end{theorem}
\section{Lecture 27 -- 4/9}
\subsection{18.1 Fields of Fractions}
\textbf{Ordered-pair set and equivalence relation:} Let $D$ be an integral domain. Define \[S=\{(a,b):a,b\in D,b\neq0\}\] Declare a relation $\sim$ on $S$ by \[(a,b)\sim(c,d)\iff ad=bc.\]

\textbf{Lemma:}  The relation $\sim$ on $S$ is an equivalence relation.

\begin{example}
  \textbf{Verifying the equivalence relation in $\Z$:} Work in $D=\Z$. Check that $(2,4)\sim(3,6)$ and $(3,6)\sim(5,10)$, and hence $(2,4)\sim(5,10)$. All three pairs represent $\frac{1}{2}$.
\end{example}

\textbf{Operations on equivalence classes:} Let $D$ be an integral domian. Denote the equivalence class of $(a,b)\in S$ under $\sim$ by $[a,b]$, and let $F_D$ be the set of all such equivalence classes. Define addition and multiplication on $F_D$ by \[[a,b]+[c,d]=[ad+bc,bd],\;[a,b]\cdot[c,d]=[ac,bd].\]

\textbf{Lemma:} The set $F_D$, together with the operations of addition and multiplication defined above, is a field.

\begin{theorem}
  Let $D$ be an integral domain. Then $D$ can be embedded in a field of fractions $F_D$, where every element of $F_D$ can be expressed as the quotient of two elements of $D$. Furthermore, $F_D$ is unique in the sense that if $E$ is any field containing $D$, then there exists an injective ring homomorphism $\psi:F_D\to E$ satistfying $\psi(a)=a$ for all $a\in D$.
\end{theorem}

\textbf{Corollary:} Let $F$ be a field of characteristic zero. Then $F$ contains a subfield isomorphic to $\Q$.

\textbf{Corollary:} Let $F$ be a field of characteristic $p$ (prime). Then $F$ contains a subfield isomorphic to $\Z_p$.

\section{Lecture 28 -- 4/16}
\subsection{20.1 Vector Spaces: Definitions and Examples}
A \textbf{vector space} $V$ over a field $F$ is a nonempty set $V$ together with two binary operations: vector addition $+V\times V\to V$ and scalar multiplication $\cdot\times V\to V$, satisfying the following axioms for all $\vec{u},\vec{v},\vec{w}\in V$ and $\alpha,\beta\in F$:

\begin{itemize}
  \item $\vec{u}+\vec{v}\in V$ (closure under addition)
  \item $\vec{u}+\vec{v}=\vec{v}+\vec{u}$ (commutativity)
  \item $(\vec{u}+\vec{v})+\vec{w}=\vec{u}+(\vec{v}+\vec{w})$ (associativity of addition)
  \item There exists a vector $\vec{0}\in V$ such that $\vec{u}+\vec{0}=\vec{u}$ for all $\vec{u}\in V$ (additive identity)
  \item For every $\vec{u}\in V$, there exists $-\vec{u}\in V$ such that $\vec{u}+(-\vec{u})=\vec{0}$ (additive inverses)
  \item $\alpha\vec{u}\in V$ (closure under scalar multiplication)
  \item $\alpha(\vec{u}+\vec{v})=\alpha\vec{u}+\alpha\vec{v}$ (distributivity over vector addition)
  \item $(\alpha+\beta)\vec{u}=\alpha\vec{u}+\beta\vec{u}$ (distributivity over scalar addition)
  \item $\alpha(\beta\vec{u})=(\alpha\beta)\vec{u}$ (associativity of scalar multiplication)
  \item There exists $1\in F$ such that $1\cdot\vec{u}=\vec{u}$ for all $\vec{u}\in V$ (multiplicative identity)
\end{itemize}

\begin{example}
  The set $\R^n$ of all $n$-tuples of real numbers, $\R^n=\{(x_1,x_2,\dots,x_n):x_i\in\R\},$ with componentwise addition and scalar multiplication is a vector space over $\R$.
\end{example}

\begin{example}
  The set $M_{m\times n}(F)$ of all $m\times n$ matrices with entries in a field $F$, with the usual matrix addition and scalar multiplication, is a vector space over $F$.
\end{example}

\begin{theorem}
  Let $V$ be a vector space over a field $F$. Then each of the following holds:
  \begin{itemize}
    \item $0\cdot\vec{v}=\vec{0}$ for every $\vec{v}\in V$, where $0\in F$ is the zero scalar and $\vec{0}\in V$ is the zero vector.
    \item $\alpha\cdot\vec{0}=\vec{0}$ for every $\alpha\in F$.
    \item $(-1)\vec{v}=-\vec{v}$ for every $\vec{v}\in V$.
    \item If $\alpha\vec{v}=\vec{0}$, then either $\alpha=0$ or $\vec{v}=\vec{0}$.
  \end{itemize}
\end{theorem}

\subsection{20.2 Subspaces}
\begin{definition}
  Let $V$ be a vector space over a field $F$ and let $W$ be a nonempty subset of $V$. We say that $W$ is a \textbf{subspace} of $V$ if $W$ is itself a vector space under the operations of $V$; equivalently, $W$ is a subspace of $V$ if and only if it is closed under vector addition and scalar multiplication:
  \begin{itemize}
    \item If $\vec{u},\vec{v}\in W,$ then $\vec{u}+\vec{v}\in W$.
    \item If $\vec{v}\in W$ and $\alpha\in F$, then $\alpha\vec{v}\in W$.
  \end{itemize}
\end{definition}

\begin{definition}
  Let $V$ be a vector space over a field $F$ and let $v_1,v_2\dots,v_n\in V$, $\alpha_1,\alpha_2,\dots,\alpha_n\in F$. Any vector of the form $w=\alpha_1v_1+\alpha_2v_2+\cdots+\alpha_nv_n$ is called a \textbf{linear combination} of $v_1,v_2,\dots,v_n$. The set of all linear combinations of a subset $S=\{v_1,\dots,v_n\}$ of $V$ is called the \textbf{span} of $S$ and is denoted by $span(S)$.
\end{definition}

\begin{theorem}
  Let $S=\{v_1,v_2,\dots,v_n\}$ be a set of vectors in a vector space $V$. Then $span(S)$ is a subspace of $V$.
\end{theorem}

\subsection{20.3 Linear Independence}
\begin{definition}
  Let $S=\{v_1,v_2,\dots,v_n\}$ be a set of vectors in a vector space $V$ over a field $F$. We say that $S$ is \textbf{linearly dependent} if there exist scalars $\alpha_1,\alpha_2,\dots,\alpha_n\in F$, not all zero, such that $\alpha_1v_1+\alpha_2v_2+\cdots+\alpha_nv_n=\vec{0}$. If the only solution to this equation is $\alpha_1=\alpha_2=\cdots=\alpha_n=0$, then $S$ is called \textbf{linearly independent}.
\end{definition}

A set $\{v_1,v_2,\dots,v_n\}$ of vectors in a vector space $V$ is linearly dependent if and only if at least one of the vectors $v_k$ can be written as a linear combination of the remaining vectors.

\begin{definition}
  A set $\{v_1,v_2,\dots,v_n\}$ of vectors in a vector space $V$ is called a $\textbf{basis}$ for $V$ if
  \begin{itemize}
    \item $\{v_1,v_2,\dots,v_n\}$ is linearly independent, and
    \item $span\{v_1,v_2,\dots,v_n\}=V$.
  \end{itemize}
\end{definition}
\begin{theorem}
  Let $\{v_1,v_2,\dots,v_n\}$ be a basis for a vector space $V$. Then every vector $v\in V$ can be written uniquely as a linear combination $v=\alpha_1v_1+\alpha_2v_2+\cdots+\alpha_nv_n$.
\end{theorem}
\begin{theorem}
  Let $V$ be a vector space that has a basis with $n$ elements. Then any two bases of $V$ have exactly $n$ elements.
\end{theorem}
\begin{definition}
  If $\{e_1,e_2,\dots,e_n\}$ is a basis for a vector space $V$, then we say the \textbf{dimension} of $V$ is $n$ and write $dim\;V=n$. If $V$ is the trivial vector space $\{\vec{0}\}$, we define $dim\;V=0$. A vector space is called \textbf{finite-dimensional} if it has a finite basis, and \textbf{infinite-dimensional} otherwise.
\end{definition}

\section{Lecture 29 -- 4/21}
\subsection{21.1 Extension Fields}
A field $E$ is an \textbf{extension field} of a field $F$ if $F$ is a subfield of $E$. The field $F$ is called the \textbf{base field}. We write $F\subset E$.

\begin{theorem}
  Let $F$ be a field and let $p(x)$ be a nonconstant polynomial in $F[x]$. Then there exists an extension field $E$ of $F$ and an element $\alpha\in E$ such that $p(\alpha)=0$.
\end{theorem}
An element $\alpha$ in an extension field $E$ over $F$ is \textbf{algebraic} over $F$ if $f(\alpha)=0$ for some nonzero polynomial $f(x)\in F[x]$.

We say the extension field $E$ is a finite extension of degree $n$ over $F$ if $E\cong F[x]/\langle p(x)\rangle$ for some irreducible polynomial $p(x)$ of degree $n$, denoted $[E:F]=n$.

An element in $E$ that is not algebraic over $F$ is \textbf{transcendental} over $F$.

\subsection{21.2 Splitting Fields}
\begin{definition}
  Let $F$ be a field and $p(x)=a_0+a_1 x+\cdots+a_nx^n$ be a nonconstant polynomial in $F[x]$. An extension field $E$ of $F$ is a \textbf{splitting field} of $p(x)$ if there exist elements $a_1,\dots,a_n$ in $E$ such that $E=F(\alpha_1,\dots,\alpha_n)$ and $p(x)=(x-\alpha_1)(x-\alpha_2)\cdots(x-\alpha_n)$.
\end{definition}

A polynomial $p(x)\in F[x]$ \textbf{splits} in $E$ if it is the product of linear factors in $E[x]$.

\begin{theorem}
  Let $p(x)\in F[x]$ be a nonconstant polynomial. Then there exists a splitting field $E$ for $p(x)$.
\end{theorem}
\begin{theorem}
  Let $p(x)$ be a polynomial in $F[x]$. Then there exists a splitting field $K$ of $p(x)$ that is unique up to isomorphism.
\end{theorem}

\section{Lecture 30 -- 4/23}
\subsection{21.3 Geometric Constructions}
A real number $\alpha$ is \textbf{constructible} if we can construct a line segment of length $|\alpha|$ in a finite number of steps from a segment of unit length by using a straightedge and compass.

\begin{theorem}
  The set of all constructible real numbers forms a subfield $F$ of the field of real numbers.
\end{theorem}

\textbf{Lemma}: If $\alpha$ is a constructible number, then $\sqrt{\alpha}$ is a constructible number.

\textbf{Lemma}: Let $F$ be a field of constructible numbers. Then the points determined by the intersections of lines and circles in $F$ lie in the field $F(\sqrt{\alpha})$ for some $\alpha$ in $F$.

\begin{theorem}
  A real number $\alpha$ is a constructible number if and only if there exists a sequence of fields \[\Q=F_0\subset F_1\subset F_2\subset\cdots\subset F_k\] such that $F_i=F_{i-1}(\sqrt{\alpha_i})$ with $\alpha_i\in F_i$ and $\alpha\in F_k$. In particular, there exists an integer $k>0$ such that $[\Q(\alpha):\Q]=2^k$.
\end{theorem}

\section{Lecture 31 -- 4/28}
\subsection{22.1 Structure of Finite Fields}
\begin{definition}
  A field $F$ has \textbf{characteristic} $p$ if $p$ is the smallest positive integer such that $p\alpha=0$ for every nonzero element $\alpha\in F$. If no such integer exists, $F$ has characteristic $0$. It is known that $p$ must be prime whenever the characteristic is finite.
\end{definition}
\begin{theorem}
  If $F$ is a finite field, then the characteristic of $F$ is $p$ for some prime $p$.
\end{theorem}
\begin{theorem}
  If $F$ is a finite field of characteristic $p$, then $|F|=p^n$ for some $n\in\N$.
\end{theorem}
\begin{definition}
  Let $F$ be a field. A polynomial $f(x)\in F[x]$ of degree $n$ is \textbf{separable} if it has $n$ distinct roots in its splitting field, i.e. it factors into distinct linear factors over that splitting field. An extension $E$ of $F$ is a \textbf{separable extension} of $F$ if every element of $E$ is the root of some separable polynomial in $F[x]$.
\end{definition}
\begin{definition}
  Let $f(x)=a_0+a_1x+\cdots+a_nx^n\in F[x]$. The \textbf{derivative} of $f(x)$ is   $f'(x)=a_1+2a_2x+\cdots+na_nx^{n-1}$. This is a formal derivative, and it is not defined in terms of limits as in calculus.
\end{definition}
\begin{theorem}
  Let $F$ be a field and $f(x)\in F[x]$. Then $f(x)$ is separable if and only if $gcd(f(x),f'(x))=1$.
\end{theorem}
\begin{theorem}
  Let $p$ be a prime and $D$ an integral domain of characteristic $p$. Then for all $a,b\in D$ and all positive integers $n$, $(a+b)^{p^n}=a^{p^n}+b^{p^n}$.
\end{theorem}
\begin{theorem}
  For every prime $p$ and every positive integer $n$, there exists a finite field with $p^n$ elements. Furthermore, any field of order $p^n$ is isomorphic to the splitting field of $x^{p^n}-x$ over $\Z_p$.
\end{theorem}
\begin{definition}
  The unique finite field with $p^n$ elements (up to isomorphism) is called the \textbf{Galois field of order $p^n$} and is denoted $GF(p^n)$.
\end{definition}
\begin{theorem}
  Every subfield of $GF(p^n)$ has order $p^m$ where $m\mid n$. Conversely, for each $m$ with $m\mid n$ there exists a unique subfield of $GF(p^n)$ isomorphic to $GF(p^m)$.
\end{theorem}

\end{document}
